\section{Miller sequences and essential incompleteness}\label{sec:miller}

For the sequence $a$ constructed by Miller
\cite[Theorem~6.3 and the following paragraph]{Miller2004}, the index
sets $\{n:a(n)<1/2\}$ and $\{n:a(n)>1/2\}$ are separable relative to
each binary representation of $a$. No single separator is computable
uniformly from all representations. We use the formulation of Bauer
and Hanson \cite[Section~2.1 and Theorem~2.2]{BauerHanson} to derive
incompleteness for an extension of the theory of real closed ordered
fields by rational bounds on the coordinates of $a$.

Let $2^{\N}=\{0,1\}^{\N}$ be Cantor space, the set of infinite binary
sequences with the product topology of the discrete space $\{0,1\}$.
Fix a recursive bijection $\langle-,-\rangle:\N^2\to\N$ and an
acceptable enumeration $(\Phi_e^p)_{e\in\N}$ of partial numerical
oracle computations, where $p\in2^{\N}$. Tuple inputs and outputs use
fixed recursive codings. Let $(q_j)_{j\in\N}$ be a recursive enumeration
of the rationals. An oracle program with index $d\in\N$
\emph{computes a real $r$ from $p$}
if $\Phi_d^p(k)$ is defined and
$|q_{\Phi_d^p(k)}-r|<2^{-k}$ for every $k\in\N$.
Give $Q=[0,1]^{\N}$ its product topology and define
\[
 \rho:2^{\N}\longrightarrow Q,\qquad
 \rho(p)(n)=\sum_{j\geq0}p(\langle n,j\rangle)2^{-j-1}.
\]
For each $n$, the subsequence $j\mapsto p(\langle n,j\rangle)$
is a binary expansion of $\rho(p)(n)$.
A \emph{binary representation} of $z\in Q$ is a sequence
$p\in2^{\N}$ with $\rho(p)=z$. Write
$P_z=\rho^{-1}(\{z\})$ for the set of these representations.
The map $\rho$ is continuous and onto, so every $P_z$ is nonempty
and compact.

Miller's construction gives a sequence $a\in Q$ such that, for every
$d\in\N$ and $r\in[0,1]$,
\begin{equation}\label{eq:miller-property}
 \left[
  (\forall p\in P_a)(\forall k\in\N)(\exists j\in\N)\;
  \bigl(\Phi_d^p(k)\mathord\downarrow=j
       \ \land\ |q_j-r|<2^{-k}\bigr)
 \right]\ \Longrightarrow\ a(d)=r.
\end{equation}
Such an $a$ is called a \emph{Miller sequence}. The program index $d$
is fixed. The rational approximation $q_j$ may depend on both the
representation $p$ and the precision $k$.

\subsection{Computation independent of the binary representation}
\label{sec:miller-category}

\begin{lemma}\label{lem:miller-unanimity}
There is one oracle program which, from $p\in P_z$ and numerical inputs
$e,n$, computes the partial function $\tau_z:\N^2\rightharpoonup\N$
defined by
\[
 \tau_z(e,n)=m
 \quad\Longleftrightarrow\quad
 (\forall q\in P_z)\;\Phi_e^q(n)\mathord\downarrow=m.
\]
Its value and its domain of definition depend only on $z$, not on the
chosen binary representation $p$.
\end{lemma}
\begin{proof}
For a finite binary sequence $\sigma$, let
$[\sigma]=\{p\in2^{\N}:p\mathbin{\upharpoonright}|\sigma|=\sigma\}$
be its cylinder in Cantor space.
The set $\rho([\sigma])$ is a finite union of closed dyadic boxes
in $Q$, computable from $\sigma$, with restrictions on finitely
many coordinates. For each affected coordinate $n$, complete the
prescribed binary digits to a finite initial segment of
$j\mapsto p(\langle n,j\rangle)$. There are only finitely many such
completions.
Given $p\in P_z$, put
\[
 b_{n,N}(p)=\sum_{j<N}p(\langle n,j\rangle)2^{-j-1}.
\]
Then $b_{n,N}(p)\le z(n)\le b_{n,N}(p)+2^{-N}$. For $c\in\mathbb Q$,
\[
 \begin{aligned}
 c<z(n)&\iff(\exists N)\ c<b_{n,N}(p),\\
 z(n)<c&\iff(\exists N)\ b_{n,N}(p)+2^{-N}<c .
 \end{aligned}
\]
These tests use finitely many bits for each $N$ and work for every
$p\in P_z$, including either expansion of a dyadic coordinate.
A point lies outside a closed box if it violates one of its interval
bounds. For each $\sigma$, wait until a violated bound has been found
for every box in $\rho([\sigma])$. Dovetailing over $\sigma$ enumerates
the cylinders $[\sigma]$ for which $z\notin\rho([\sigma])$.
Their union is exactly
$2^{\N}\setminus P_z$: if $\rho(q)\ne z$, a sufficiently long initial
segment of $q$ already forces a different value in one coordinate.

For fixed $e,n,m$, let
$H_{e,n,m}=\{q\in2^{\N}:\Phi_e^q(n)\mathord\downarrow=m\}$.
Finite halting computations enumerate this open set as a union of
cylinders, uniformly in $e,n,m$. For each $m$, search for a finite cover
of Cantor space by cylinders from $H_{e,n,m}$ and
$2^{\N}\setminus P_z$, dovetailing these searches over $m$. Whether finitely
many cylinders cover Cantor space is decidable. For a nonempty list,
check whether every word of the maximum listed length extends a listed
prefix. Such a cover exists
exactly when $P_z\subseteq H_{e,n,m}$, by compactness. Since $P_z$ is
nonempty, at most one $m$ satisfies this condition. Thus the search
computes $\tau_z$ from every binary representation of $z$.
\end{proof}

The objects of $\mathcal E_a$ are generated from $\N$ by binary
Cartesian products and disjoint unions. Fix for each object $X$ a
recursive bijection $\iota_X:X\to\N$. Take $\iota_{\N}=1_{\N}$
and $\iota_{X\times Y}(x,y)=\langle\iota_X(x),\iota_Y(y)\rangle$.
For $X+Y$, code the left summand by $2\iota_X(x)$ and the right
summand by $2\iota_Y(y)+1$.
For a partial function
$f:X\rightharpoonup Y$, write $\widehat f=\iota_Y f\iota_X^{-1}$.
Then
\[
 f\in\mathcal E_a(X,Y)
 \quad\Longleftrightarrow\quad
 (\exists e\in\N)(\forall p\in P_a)(\forall n\in\N)\;
       \Phi_e^p(n)\simeq\widehat f(n).
\]
Here $u\simeq v$ means that both sides are undefined or both are
defined with the same value. Thus a single index must compute the
entire partial function from each representation of $a$. Two morphisms
are equal when their partial functions are equal.

\begin{proposition}\label{prop:miller-category}
The category $\mathcal E_a$ is a locally connected recursion category
with ranges, binary restriction coproducts, and finite compatible joins.
It satisfies weak choice, and $\tau_a:\N\times\N\rightharpoonup\N$
is a Turing morphism.
\end{proposition}
\begin{proof}
Composition of the representing programs gives the same composite
partial function for every binary representation of $a$.
The recursive projections, pairings,
coproduct injections, and functions defined separately on the two
summands give the restriction
products and coproducts. Restriction is computed by running a program for $f(x)$ and,
if it halts, returning $x$. To compute the range partial identity at
$y$, dovetail the computations $f(x)$ over all coded inputs $x$ and
return $y$ when one returns $y$. Although the discovered witness may
depend on the representation, the returned value and the existence of a witness
do not. Thus $\ran(f)$ belongs to $\mathcal E_a$; its leastness and
agreement law follow from those of set-theoretic images.
Dovetailing compatible computations gives their union;
compatibility ensures that the output is independent of which
computation finishes first. These operations also give zero morphisms and
the prodominical law $f\times0=0$.

For every input $x$, the singleton partial identity $1_{\{x\}}$ is
recursive. It is connected, and these partial identities detect
inequality of partial functions, so the category is locally connected.
A partial monomorphism here is injective on its domain of definition: a failure of
injectivity would be detected by two recursive constant morphisms. For an
injective morphism $f$, the same preimage search returns the unique $x$
with $f(x)=y$. Its output is consequently independent of the representation and
gives the partial inverse. This proves weak choice.

Lemma~\ref{lem:miller-unanimity} puts $\tau_a$ in $\mathcal E_a$.
Let $f:X\times\N\rightharpoonup\N$ have index $d$. The parameter
theorem for oracle programs gives a total recursive map $s_d:\N\to\N$
with
\[
 \Phi_{s_d(u)}^p(n)\simeq\Phi_d^p(\langle u,n\rangle)
 \qquad(u,n\in\N,\ p\in2^{\N}).
\]
Put $c=s_d\iota_X:X\to\N$. For $p\in P_a$ the specialized
program at $u=\iota_X(x)$ computes the partial function $f(x,-)$.
If $f(x,n)=m$, every $p\in P_a$ gives this same output; if $f(x,n)$ is
undefined, every $p\in P_a$ gives an undefined computation. The definition
of $\tau_a$ gives $f=\tau_a(c\times1_{\N})$.
Finally, the chosen recursive codings make the category isotypical.
\end{proof}

For $p\in P_a$, let $\mathrm{Rec}^{p}$ be the recursion category on
the same numerical objects, with partial $p$-recursive functions as
morphisms. The identity interpretation of partial functions defines a
faithful range functor $J_p:\mathcal E_a\to\mathrm{Rec}^{p}$.
It also preserves restriction products and finite compatible joins.
Write $\mathcal K_p=\operatorname{Split}(\mathrm{Rec}^{p})$ and let
$I_p:\mathrm{Rec}^{p}\to\mathcal K_p$ be the full inclusion.
This inclusion is a range functor.
A presentation of $\mathcal K_p$ has $p$-c.e. subsets of the
numerical objects as objects and partial $p$-recursive functions between
these carriers as morphisms. A restriction idempotent splits through its
$p$-c.e. domain of definition.
Ranges are the partial identities on images, and inverse images along
total morphisms preserve unions of domains of definition.

If $c,c':X\rightharpoonup X$ are complementary restriction
idempotents in $\mathcal E_a$, Propositions~\ref{prop:miller-category}
and~\ref{prop:disjoint-joins} give a total morphism
$\ell_c:X\to X+X$ with $\ell_cc=i_1c$ and $\ell_cc'=i_0c'$,
where $i_0,i_1:X\to X+X$ are the coproduct injections.
Since $(I_pJ_p)(1_{\N})\ne0$, each composite
$I_pJ_p$ satisfies the hypotheses of Corollary~\ref{cor:source-conditionals}.
For $X=\N$, define the recursive bijection $b:\N\to\N+\N$ by
$b(2n)=i_0(n)$ and $b(2n+1)=i_1(n)$. Applying that corollary with $b$ gives inseparable
restriction idempotents that are the partial identities on the sets
\begin{equation}\label{eq:miller-diagonal-pair}
 A_i=\{e:\tau_a(e,e)\mathord\downarrow=m
                   \text{ for some }m\equiv i\pmod2\},
 \qquad i=0,1.
\end{equation}

\subsection{The index sets below and above \texorpdfstring{$1/2$}{1/2}}
\label{sec:miller-cuts}

Define $L_a=\{n:a(n)<\tfrac12\}$ and
$R_a=\{n:a(n)>\tfrac12\}$.
Their partial identities belong to $\mathcal E_a$ by the two
finite-prefix tests in Lemma~\ref{lem:miller-unanimity}, with $c=1/2$.

\begin{lemma}\label{lem:miller-cut-reduction}
There is a total recursive map $h:\N\to\N$ such that
\[
 h^{-1}(L_a)=A_0,\qquad h^{-1}(R_a)=A_1.
\]
\end{lemma}
\begin{proof}
For $e,m,t\in\N$, let $H_{e,m,t}$ be the finite union of halting
cylinders enumerated by stage $t$ for the computation
$\Phi_e^q(e)\mathord\downarrow=m$. Choose these finite unions increasing
in $t$, and put $K_{e,m,t}=\rho(2^{\N}\setminus H_{e,m,t})\subseteq Q$.
Each $K_{e,m,t}$ is a finite union of closed dyadic boxes computable
from $e,m,t$. Compactness of $P_z$ gives
\begin{align*}
 U_e^i
 &:=\{z\in Q:\tau_z(e,e)\mathord\downarrow=m
                      \text{ for some }m\equiv i\pmod2\}\\
 &=\bigcup_{\substack{m\equiv i\pmod2\\t\in\N}}
       (Q\setminus K_{e,m,t}).
\end{align*}
In particular, $U_e^0$ and $U_e^1$ are disjoint effectively open sets.

The product topology on $Q$ is induced by the metric
\[
 d:Q\times Q\to[0,1],\qquad
 d(z,w)=\sum_{n\geq0}2^{-n-1}|z(n)-w(n)|.
\]
For a nonempty box restricting finitely many coordinates to closed
rational intervals $I_n$, the distance from $z$ to the box is the
finite sum
$\sum_n2^{-n-1}\operatorname{dist}(z(n),I_n)$ over those coordinates.
Distance to a nonempty finite union is the minimum of these distances.
Define $\delta_{e,m,t}:Q\to[0,1]$ by
$\delta_{e,m,t}(z)=d(z,K_{e,m,t})$ when $K_{e,m,t}$ is nonempty, and
$\delta_{e,m,t}(z)=1$ when it is empty. In each box formula, replacing
$z(n)$ by $b_{n,N}(p)$ changes the distance by at most
$2^{-N}\sum_n2^{-n-1}=2^{-N}$, since distance to an interval is
$1$-Lipschitz. Taking a finite minimum preserves this error bound.
Thus the distances are computable uniformly in $e,m,t$ from each
$p\in P_z$. They are continuous, and $\delta_{e,m,t}(z)>0$ exactly
when $z\notin K_{e,m,t}$.
For each $i$, take a fixed recursive enumeration
$(m_j,t_j)_{j\geq0}$ of all pairs with $m_j\equiv i\pmod2$, and define
\[
 g_e^i:Q\to[0,1],\qquad
 g_e^i(z)=\sum_{j\geq0}2^{-j-1}\delta_{e,m_j,t_j}(z).
\]
Since $0\le\delta_{e,m,t}\le1$, the tail from $j=N$ onward is at most
$2^{-N}$. Approximating each of the first $N$ distances to within $2^{-N}$
gives a rational approximation to $g_e^i(z)$ with error at most
$2^{1-N}$. This proves computability from every $p\in P_z$, and uniform
convergence gives continuity. The sum belongs to $[0,1]$ and is positive
exactly on $U_e^i$. The computable continuous function
\[
 r_e:Q\to[1/4,3/4],\qquad
 r_e(z)=\frac12+\frac{g_e^1(z)-g_e^0(z)}4
\]
is below $1/2$ exactly on $U_e^0$, and
is above $1/2$ exactly on $U_e^1$.

Given $e,k$ and $p\in2^{\N}$, take $N=k+1$ in the two
approximations to $g_e^i(\rho(p))$. Substituting them in the formula for
$r_e$ gives a rational with error at most $2(2^{1-N})/4=2^{-N}$.
Search the rational enumeration for its index. This gives one oracle
program $\Psi$ with
\[
 \bigl|q_{\Psi^p(e,k)}-r_e(\rho(p))\bigr|<2^{-k}
 \qquad(e,k\in\N,\ p\in2^{\N}).
\]
The oracle parameter theorem gives a total recursive $h$ such that
$\Phi_{h(e)}^p(k)=\Psi^p(e,k)$. For $p\in P_a$ these computations
represent the same real $r_e(a)$, so the Miller property gives
$a(h(e))=r_e(a)$.
Together with $a\in U_e^i\Longleftrightarrow e\in A_i$, this proves
the asserted inverse-image equalities.
\end{proof}

\begin{corollary}\label{cor:miller-cuts-inseparable}
The restriction idempotents $1_{L_a},1_{R_a}:\N\rightharpoonup\N$ are inseparable in
$\mathcal E_a$. Every $p\in P_a$ computes a separator
of the underlying sets $L_a,R_a$.
\end{corollary}
\begin{proof}
If complementary restriction idempotents in $\mathcal E_a$ separated
$1_{L_a},1_{R_a}$, their inverse images under $h$ would separate
$1_{A_0},1_{A_1}$, contrary to
Corollary~\ref{cor:source-conditionals} as applied above.
For $p\in P_a$, put $C_p=\{n:p(\langle n,0\rangle)=0\}$. A binary
expansion beginning with $0$ represents a number at most $1/2$,
and one beginning with $1$ represents a number at least $1/2$. Hence
\[
 L_a\subseteq C_p\subseteq\N\setminus R_a
 \qquad(p\in P_a).
\]
One program decides $C_p$ by reading the first digit. At $a(n)=1/2$
the two expansions $0.0111\ldots$ and $0.1000\ldots$ give different
answers. If a single set $C$ separating the cuts were decided by one
program for every $p\in P_a$, the complementary partial identities
$1_C$ and $1_{\N\setminus C}$ would separate the cuts in $\mathcal E_a$.

\end{proof}

If $J_p\tau_a$ were a Turing morphism, Corollary~\ref{cor:source-conditionals},
applied to $I_p:\mathrm{Rec}^{p}\to\mathcal K_p$ and its
$p$-recursive conditionals associated with complementary restriction
idempotents, would make $1_{A_0},1_{A_1}$ inseparable in
$\mathrm{Rec}^{p}$. The $p$-decidable set $h^{-1}(C_p)$ separates
$A_0,A_1$. Hence $J_p\tau_a$ is not a Turing
morphism, although $\mathrm{Rec}^{p}$ has a Turing morphism.


\subsection{A theory of real closed fields}
\label{def:miller-splitting}

A subset $D\subseteq X$ is \emph{$\mathcal E_a$-semidecidable} if
its partial identity $1_D$ belongs to $\mathcal E_a$. In the chosen
coding this says
\[
 (\exists e\in\N)(\forall p\in P_a)\quad
 \operatorname{dom}(\Phi_e^p)=\iota_X(D).
\]
Indeed, a computation with this domain computes the partial identity
if it returns its input whenever it halts.

We use the presentation of $\operatorname{Split}(\mathcal E_a)$
in which such a subset $D\subseteq X$ represents $(X,1_D)$.
For another $\mathcal E_a$-semidecidable subset $E\subseteq Y$,
a morphism $D\rightharpoonup E$ is represented by
a morphism $f:X\rightharpoonup Y$ of $\mathcal E_a$ satisfying
$f=1_Ef1_D$; it is total when $\bar f=1_D$.
Put $\mathcal D_a=\operatorname{Tot}(\operatorname{Split}(\mathcal E_a))$.

\begin{proposition}\label{prop:miller-theory-realization}
The category $\mathcal D_a$ is regular, with finite coproducts that
are disjoint and stable under pullback. Every monomorphism of
$\mathcal D_a$ has a partial inverse in
$\operatorname{Split}(\mathcal E_a)$, and there is a restriction equivalence
\[
 \operatorname{Par}(\mathcal D_a,\operatorname{Mono})
       \simeq\operatorname{Split}(\mathcal E_a).
\]
For $V=\N$ and a recursive signature of bounded finite arity, the
syntax and list algebras of Definition~\ref{def:categorical-syntax}
exist in $\mathcal D_a$ and are stable under parameters.
For every axiom subobject of the sentence object, the inference
endofunctor of Theorem~\ref{thm:first-order-theory} has an initial
algebra of finite derivations.
\end{proposition}
\begin{proof}
A singleton is terminal. Products have carriers $C\times D$, and the
equalizer of $f,g:C\to D$ has carrier $\{x\in C:f(x)=g(x)\}$.
Parallel semidecision and numerical equality semidecide these carriers.
The empty carrier and disjoint unions give finite coproducts.
For $f:C\to D+E$, write $i_D,i_E$ for the coproduct injections and put
\[
 C_D=f^{-1}(i_D(D)),\qquad C_E=f^{-1}(i_E(E)).
\]
Each carrier is semidecidable by running $f$ and reading its summand.
They are disjoint. The inclusions induce an isomorphism
$C_D+C_E\to C$ whose inverse sends $x$ to the summand determined by
$f(x)$. Thus pulling back the coproduct injections along $f$ gives a
coproduct decomposition of $C$.

For $f:C\to D$, dovetail the computations of $f(x)$ over all coded
inputs $x$. This semidecides $f(C)$ and gives the image factorization
$C\twoheadrightarrow f(C)\hookrightarrow D$. For a surjection
$e:C\to D$, the kernel pair has carrier
\[
 C\times_D C=\{(x,x')\in C\times C:e(x)=e(x')\}.
\]
Let $p_1,p_2:C\times_D C\to C$ be its projections. If $h:C\to Z$
satisfies $hp_1=hp_2$, then there is a unique $k:D\to Z$ with
\[
 k(y)=h(x)\quad\text{whenever } e(x)=y.
\]
Indeed, on input $y$, dovetail the computations of $e(x)$ until one
returns $y$, then return $h(x)$. The search succeeds exactly on $D$.
Different binary representations of $a$ may yield different witnesses
$x$, but $hp_1=hp_2$ makes the returned value independent of the
witness. Thus $k$ is a morphism of $\mathcal D_a$, and $e$ is the coequalizer
of its kernel pair. Conversely, if $f$ is a regular epimorphism, its
extremality makes the inclusion $f(C)\hookrightarrow D$ an
isomorphism, so $f$ is surjective. For any $g:B\to D$, the projection
$B\times_D C\to B$ is surjective when $e$ is. Hence regular
epimorphisms are stable under pullback, proving regularity.

Singleton morphisms detect failure of injectivity, so the monomorphisms
of $\mathcal D_a$ are the injections. For an injection $m:C\to D$,
the same preimage search computes the partial morphism
$m^\dagger:D\rightharpoonup C$ with
\[
 m^\dagger m=1_C,\qquad mm^\dagger=1_{m(C)}.
\]
For $y\in m(C)$ the witness is unique. The span
$C\xleftarrow{m}B\xrightarrow{g}D$ represents the partial morphism
$gm^\dagger:C\rightharpoonup D$. Conversely, a partial morphism
$u:C\rightharpoonup D$ is represented by
$C\hookleftarrow\operatorname{dom}(u)
\xrightarrow{u|_{\operatorname{dom}(u)}}D$.
For $v:D\rightharpoonup E$, the domain of $vu$ is the inverse image of $\operatorname{dom}(v)$
under the total restriction of $u$ to $\operatorname{dom}(u)$.
This is the pullback used to compose the spans. The two assignments
therefore give the stated restriction equivalence.

Equality on the variable object $V=\N$ is decidable. Encode terms and formulas
as finite labeled ordered trees with distinct numerical labels for
the constructors, and lists of variable names as finite sequences. Recursive parsing
and the recursive symbol carriers make these carriers and their
constructors recursive. Fix an arbitrary object $P$ of $\mathcal D_a$
and a $K_P$-algebra on two objects over $P$. On input $(w,t)$ or
$(w,\varphi)$, first semidecide $w\in P$ and check the tree, then
evaluate it from its leaves to its root using the algebra operations
in the fiber over $w$. Each operation preserves that fiber.
Induction on the tree proves termination and gives the same result
for every $p\in P_a$, because each algebra operation has a common
domain and value. Thus evaluation defines total morphisms from
$P\times\mathsf{Tm}$ and $P\times\mathsf{Fm}$. They satisfy the algebra
equations, and induction
proves uniqueness. This proves syntax initiality in every slice
over $P$. Evaluation of finite lists, retaining the same parameter
$w$, proves the corresponding assertion for the list algebra.

The free-variable morphism of Appendix~\ref{app:syntax-theories}
computes finite lists of occurrences. The formula-substitution morphism
has codomain $1+\mathsf{Fm}$, with injections
$\fail_{\mathsf{Fm}}:1\to1+\mathsf{Fm}$ and
$\ok_{\mathsf{Fm}}:\mathsf{Fm}\to1+\mathsf{Fm}$.
Its value belongs to the first summand when substitution is inadmissible
and to the second summand otherwise. Its defining recursion
compares the bound and substituted variables and uses the finite
membership tests in~\eqref{eq:quantifier-fold}. These computations
are recursive functions. The pullback in
Proposition~\ref{prop:categorical-substitution} therefore gives the
recursive free-for domain and substitution on that domain. The
sentence carrier, where the free-variable list is empty, is recursive.

Represent an axiom subobject by its semidecidable image in the
sentence carrier. Use the seven logical schemes, modus ponens,
and generalization of Appendix~\ref{app:first-order-derivations}.
A finite derivation tree records its rule instance and ordered
children at every node. Tree shape, logical instance conditions,
and agreement of child conclusions with the required premises are
recursive tests. Run these tests together with the finitely many
semidecision procedures for membership of nonlogical axiom leaves.
Together these procedures semidecide a carrier
$P_A$ with the same membership outcome for every $p\in P_a$.
Reading the root gives a total conclusion morphism
$c_A:P_A\to\mathsf{Fm}$, and adjoining a labeled root gives an
$H_A$-algebra.

For any $H_A$-algebra, first semidecide that the input is a valid
derivation, then evaluate it from its leaves to its root. Induction
shows that each intermediate value has the prescribed conclusion,
so the next algebra operation is applied on its domain. The
evaluation terminates on every valid derivation. The operations'
common domains and values make its domain and result independent
of $p\in P_a$, so it is a morphism of $\mathcal D_a$ over
$\mathsf{Fm}$. Its recursive equations make it an algebra morphism;
induction on derivations proves uniqueness. Thus the derivation
algebra is initial.
\end{proof}

\begin{lemma}\label{lem:miller-parameterized-deduction}
Let $X$ be an object of $\mathcal D_a$. The first-order language
with variables $\N$, equality, and a nullary predicate for each
label in $X$ has the syntax and list algebras of
Definition~\ref{def:categorical-syntax}, stable under parameters.
For this language and for every recursive signature covered by
Proposition~\ref{prop:miller-theory-realization}, finite derivation
algebras are also stable under parameters. More precisely, if $P$
is an object of $\mathcal D_a$, $\mathsf{Fm}$ is the formula object,
$S$ is the sentence object, and
$A\hookrightarrow P\times S$ is an axiom subobject over $P$, the
corresponding inference endofunctor on
$\mathcal D_a/(P\times\mathsf{Fm})$ has an initial algebra.
\end{lemma}
\begin{proof}
Represent $X$ and $P$ by $\mathcal E_a$-semidecidable numerical
carriers. In the nullary language the term and formula functor of
Definition~\ref{def:categorical-syntax} reduces to
\[
 K(Y,Z)=(\N,\;Y^2+X+Z+Z^2+\N\times Z).
\]
Its formula constructors are equality, a predicate with label in $X$,
negation, implication and universal quantification. To recognize a
formula tree, check its shape and run the membership tests for its
finitely many predicate labels in parallel. This gives a semidecision
uniform in $p\in P_a$. Occurrence lists and substitution are computed
on these trees as in Proposition~\ref{prop:miller-theory-realization}.

Let $\iota$ denote the constructor algebra, and let $\iota_P$ be its
pullback to the slice over $P$. For a $K_P$-algebra
$\alpha:K_P(Y,Z)\to(Y,Z)$, evaluate each constructor at the given
parameter $w\in P$. If $t$ is a valid term or formula tree, induction
on $t$ shows that evaluation terminates in the fiber over $w$ and gives
the same value for every $p\in P_a$. Testing $w\in P$ and validity
before evaluation gives a pair of total morphisms over $P$,
\[
 h:(P\times\Tm,P\times\Fm)\longrightarrow(Y,Z),
 \qquad h\iota_P=\alpha K_P(h).
\]
The equation determines $h$ on every constructor, hence on every tree.
For lists the same recursion on the empty list and on cons gives the
unique algebra morphism from $P\times\List\N$ in the slice.

For either language put $F_P=P\times\Fm$. Let $L$ be the coproduct of the seven parameter objects for the logical
axiom schemes. Let $R_1$ and $R_2$ be the objects of generalization and
modus ponens instances of Appendix~\ref{app:first-order-derivations}. For
$A\hookrightarrow P\times S$ the rule objects are
\[
 R_0^P=(P\times L)+A,\qquad R_1^P=P\times R_1,\qquad
 R_2^P=P\times R_2.
\]
Giving every premise and the conclusion of each rule the same parameter
$w$ defines the inference endofunctor $H$ on $\mathcal D_a/F_P$. A derivation at $w$ is
a finite tree of these rule instances in which the ordered child
conclusions agree with the premises at each node.
For a nonlogical axiom leaf labelled $\varphi$, the test is
$(w,\varphi)\in A$. Run all leaf tests, the syntax and rule-instance
tests, and the test $w\in P$ in parallel. Agreement of codes is
decidable. Thus the carrier $D_A$ of valid pairs $(w,t)$ is
$\mathcal E_a$-semidecidable. Reading the root defines
$c_A:D_A\to F_P$ by $c_A(w,t)=(w,\varphi)$ for root conclusion
$\varphi$. Adjoining a rule above its premises gives
$\delta:H(D_A,c_A)\to(D_A,c_A)$.

Let $\beta:H(Q,q)\to(Q,q)$ be an algebra over $F_P$. Evaluate each
valid tree from the leaves to the root using $\beta$. Inductively,
each child value lies over its premise at $w$, so the rule and its
child values belong to the domain of the corresponding operation of
$\beta$. For every $p\in P_a$ it returns the same value over the rule
conclusion. With the test for membership in $D_A$
this gives a total morphism $e:D_A\to Q$ satisfying
\[
 qe=c_A,\qquad e\delta=\beta H(e).
\]
The second equation forces the value at each leaf and at each node
once its children have been evaluated. Induction gives uniqueness,
so $(D_A,c_A,\delta)$ is initial.
\end{proof}

Consider the language of ordered fields with additional constants
$(c_n)_{n\in\N}$. Its symbol carriers are recursive and its arities
are at most two, so Proposition~\ref{prop:miller-theory-realization}
applies. Write $S$ for the object of sentence codes. Let
$\Gamma_a\subseteq S$ consist of a recursive axiomatization of
real closed ordered fields together with the sentences
\begin{equation}\label{eq:miller-theory}
 q<c_n\quad\text{if }q<a(n),\qquad
 c_n<q\quad\text{if }a(n)<q,
 \qquad(q\in\mathbb Q, n\in\N).
\end{equation}
For each $q=u/v$ with $u\in\mathbb Z$ and $v$ a positive integer,
the two bounds stand for the sentences
$\underline{u}<\underline{v}\,c_n$ and
$\underline{v}\,c_n<\underline{u}$, respectively, where $\underline{j}$
is the ring term for the integer $j$.
From any $p\in P_a$, enumerate the recursive field axioms and
dovetail the two strict-inequality tests of
Lemma~\ref{lem:miller-unanimity} over all $n,u,v$ with $v>0$.
Whenever a test succeeds, enumerate its sentence. For each $p\in P_a$
this enumerates exactly $\Gamma_a$, so $\Gamma_a$ is
$\mathcal E_a$-semidecidable. Let $T_a$ be the sentence theory
generated by $\Gamma_a$. The real field with $c_n$ interpreted as $a(n)$
is a model of $T_a$.

\begin{theorem}\label{thm:miller-incompleteness}
Let $a$ be a Miller sequence with property~\eqref{eq:miller-property}.
In the language of ordered fields with constants $c_n$ for $n\in\N$,
let $T_a$ be the deductive closure of the axioms for real closed
ordered fields together with the rational bounds in~\eqref{eq:miller-theory}.
For every consistent extension of $T_a$ in the same language with
an $\mathcal E_a$-semidecidable axiom set, the partial identities on its
theorem and refutation sets are inseparable in
$\operatorname{Split}(\mathcal E_a)$.
In particular, every such extension is incomplete.
\end{theorem}
\begin{proof}
The recursive set $S$ of sentence codes is an
object of $\operatorname{Split}(\mathcal E_a)$. Negation and insertion
of a constant into an inequality give recursive total morphisms
\[
 \nu:S\to S,\qquad
 \nu(\ulcorner\varphi\urcorner)=\ulcorner\neg\varphi\urcorner,
 \qquad
 \theta:\N\to S,\qquad
 \theta(n)=\ulcorner 2c_n<1\urcorner.
\]
For a theory $T$ with an $\mathcal E_a$-semidecidable axiom set,
let $A_T\hookrightarrow S$ be its axiom subobject. By
Proposition~\ref{prop:miller-theory-realization}, the inference
endofunctor $H_{A_T}$ has an initial derivation algebra. Pulling its conclusion morphism $\rho_T:P_T\to\mathsf{Fm}$
back along the sentence inclusion $S\hookrightarrow\mathsf{Fm}$ gives
$\mathrm{Prf}_T=P_T\times_{\mathsf{Fm}}S$ with projection
$c_T:\mathrm{Prf}_T\to S$.
Theorem~\ref{thm:first-order-theory} identifies its image with the
sentence theory generated by $A_T$.
Proposition~\ref{prop:miller-theory-realization} supplies a partial
inverse for this image monomorphism in
$\operatorname{Split}(\mathcal E_a)$.
Lemma~\ref{lem:range-regular-image} therefore identifies
$t_T=\ran(c_T)$ with the partial identity on the theorem set in
$\operatorname{Split}(\mathcal E_a)$.
The inverse image $\nu^*t_T$ is the partial identity on the set of
sentence codes whose negations are provable in $T$.

Let $T\supseteq T_a$ be a consistent extension with
$\mathcal E_a$-semidecidable axioms. Every member of $\Gamma_a$ is derivable in $T$. Replacing the uses of those
axioms in a finite derivation gives $t_{T_a}\le t_T$. For $n\in L_a$, the
sentence $2c_n<1$ belongs to $\Gamma_a$. For $n\in R_a$, the
axiom $1<2c_n$ and asymmetry of the strict order give
$\neg(2c_n<1)$.
Thus
\[
 1_{L_a}\le\theta^*t_{T_a}\le\theta^*t_T,\qquad
 1_{R_a}\le\theta^*\nu^*t_{T_a}\le\theta^*\nu^*t_T .
\]
The cuts remain inseparable in the splitting by
Corollary~\ref{cor:miller-cuts-inseparable}, since the restriction
idempotents and their order on $(\N,1_{\N})$ are those of $\mathcal E_a$.
Consistency gives $t_T\nu^*t_T=0_S$. Corollary~\ref{cor:domain-incompleteness}
applied to $\theta$ now gives inseparability of $t_T,\nu^*t_T$ and
$t_T\vee\nu^*t_T\ne1_S$. A sentence code outside this union
is the code of a sentence $\varphi$ for which
\[
 T\nvdash\varphi \qquad\text{and}\qquad T\nvdash\neg\varphi .
\]

\end{proof}

\subsection{The associated recursion categories}

Write $\mathscr C_a=\Split(\mathcal E_a)$, and let
$t_a:\mathsf T_a\hookrightarrow S$ be the inclusion representing $T_a$. A sentence theory $u:U\hookrightarrow S$ in
$\mathcal D_a$ is understood to be closed under the field-language
deduction rules above.

\begin{corollary}\label{cor:miller-associated-categories}
Let $a$ be a Miller sequence, and let $u:U\hookrightarrow S$ be a
consistent internal sentence theory extending $t_a$. The category
\[
 \mathscr C_a[U]=\Par(\mathcal D_a/U,\Mono)
\]
is a Turing category, and its internal theory
\[
 1_U\times u:U\times U\hookrightarrow U\times S
\]
is essentially incomplete. Its theorem and refutation domains are
inseparable. Let $\delta_u:S\rightharpoonup\Two$ be the partial
morphism with domain $u\vee\nu^*u$ taking the value true on $u$
and false on $\nu^*u$. Then
$1_U\times\delta_u:U\times S\rightharpoonup U\times\Two$
has no total extension in $\mathscr C_a[U]$.

If $u':U'\hookrightarrow S$ is another consistent internal
sentence theory and $u=u'i$ for $i:U\hookrightarrow U'$, then
postcomposition induces a fully faithful restriction functor
\[
 \Sigma_i:\mathscr C_a[U]\longrightarrow\mathscr C_a[U'].
\]
This functor is not essentially surjective when $i$ is proper.
\end{corollary}
\begin{proof}
Proposition~\ref{prop:miller-theory-realization} supplies regularity,
disjoint stable finite coproducts, and the partial inverses of
monomorphisms. Lemma~\ref{lem:miller-parameterized-deduction}
supplies the syntax and finite derivations with parameters,
including the languages with nullary predicate symbols used in the
general construction.

Soundness in the real-field model of $T_a$ makes $t_a$ disjoint
from its refutation subobject. The total morphism
$\theta:\N\to S$, $\theta(n)=\ulcorner 2c_n<1\urcorner$,
reduces the inseparable pair
$1_{L_a},1_{R_a}$ to these subobjects, as in the proof of
Theorem~\ref{thm:miller-incompleteness}.
Corollary~\ref{cor:domain-incompleteness} therefore applies to
every consistent internal sentence theory extending $t_a$.
In particular, $u$ is essentially incomplete and its theorem and
refutation domains are inseparable.

Let $q:1\to S$ select the code of $0=0$. Since $T_a$ proves $0=0$,
$q$ factors as $u q_U$ for $q_U:1\to U$. Suppose complementary
subobjects of $U\times S$ in the slice separated $1_U\times u$
from $1_U\times\nu^*u$. Their inverse images along
\[
 \langle q_U !_S,1_S\rangle:S\longrightarrow U\times S
\]
would be complementary subobjects of $S$ containing $u$ and
$\nu^*u$, respectively. This contradicts inseparability in the base.
Theorems~\ref{thm:theory-extension} and~\ref{thm:associated-incompleteness}
give the conclusions about the category, its internal theory, the
partial Boolean morphism, and the comparison functor.
\end{proof}

